% Preliminary Report — Sophakotra's sections
% Operation Classification + Baselines + Literature Review
%
% This is a self-contained, compilable document so you can check your sections
% in isolation. To merge into the team's shared prelim_report.tex, copy the
% \section blocks and the relevant \bibitem entries.
%
% All numbers below are the actual outputs of:
%   python -m src.data.pseudo_labeler
%   python -m src.classifier.tfidf_classifier
%   python -m src.baselines.baseline3_direct_llm

\documentclass[11pt]{article}
\usepackage[margin=1in]{geometry}
\usepackage{booktabs}
\usepackage{amsmath}
\usepackage{hyperref}

\title{Operation-Aware Biomedical Text Simplification\\
\large Preliminary Report --- Operation Classification and Baselines}
\author{Sophakotra}
\date{July 2026}
% Preliminary report - Rishabh's sections (Introduction, Data, Related Work)
% Self-contained so it compiles on its own. When merging into the team
% template, lift the \section blocks and the \bibitem entries as needed.

\documentclass[11pt]{article}
\usepackage[margin=1in]{geometry}
\usepackage{times}
\usepackage{url}

\title{Diagnose Before Simplifying:\\ Operation-Aware Biomedical Text Simplification}
\author{Sruthilaya \and Rishabh \and Sophakotra \and Zihao}
\date{}

\begin{document}
\maketitle

%======================================================================
\section{Operation Classification}
%======================================================================

The central contribution of our system is a \emph{diagnostic} step: before
simplifying a span of biomedical text, we predict \emph{which} simplification
operation it requires. We adopt the three-operation taxonomy of
Heineman et al.~\cite{salsa}: \textbf{Substitution} (replacing a jargon term
with a plain-language equivalent), \textbf{Explanation} (adding words to unpack
a concept), and \textbf{Generalization} (compressing or removing detail). While
prior PLABA systems apply a single undifferentiated simplification model, we
argue that routing each span to the appropriate operation is what preserves
safety-critical information.

\paragraph{Pseudo-labeling.}
Span-level operation labels are not available in PLABA, so we derive them by
weak supervision from the source$\rightarrow$target sentence pairs. The
\emph{type} of edit a human writer performed leaves a measurable fingerprint on
the text: an Explanation lengthens the sentence, a Generalization shortens it,
and a Substitution changes many characters while keeping the length roughly
constant. Concretely, let $r = |t| / |s|$ be the target/source word-length
ratio. We label a pair Explanation if $r \ge 1.15$, Generalization if
$r \le 0.85$, and otherwise Substitution, using normalized character-level edit
distance to confirm that similar-length pairs were genuinely rewritten rather
than copied. Edit distance is the most direct measurable signal of the kind of
change made, and requires no gold annotation --- standard weakly-supervised
practice.

\paragraph{Data and leakage control.}
The PLABA release (\texttt{data.json}) is the full corpus and contains the
abstracts held out in the official validation and test splits. To avoid
train/test leakage, we exclude every validation and test PMID from the training
pairs at the source level. This yields \textbf{6{,}307} pseudo-labeled training
sentence pairs, distributed across the three operations as shown in
Table~\ref{tab:dist}. We evaluate on \textbf{1{,}431} sentence pairs built from
the held-out validation PMIDs and pseudo-labeled identically.

\begin{table}[h]
\centering
\begin{tabular}{lrr}
\toprule
Operation & Train pairs & Share \\
\midrule
Substitution   & 2{,}723 & 43.2\% \\
Explanation    & 2{,}167 & 34.4\% \\
Generalization & 1{,}417 & 22.5\% \\
\midrule
Total          & 6{,}307 & 100\% \\
\bottomrule
\end{tabular}
\caption{Pseudo-labeled operation distribution on the leakage-free training set.}
\label{tab:dist}
\end{table}

\paragraph{Three-model progression.}
We plan a controlled progression of classifiers, each isolating one variable:
(i) \textbf{TF-IDF + Logistic Regression} tests whether surface lexical features
alone predict the operation; (ii) \textbf{BERT-base}~\cite{bert} adds general
contextual pretraining; and (iii) \textbf{BioBERT}~\cite{biobert} adds
biomedical-domain pretraining on top. Reporting all three localizes exactly
where any gains originate. This preliminary report covers the TF-IDF baseline;
the contextual models follow in the full paper.

\paragraph{Preliminary results.}
The classifier receives only the \emph{source} sentence, since in the full
pipeline the operation must be chosen before the simplified text is generated.
Using $(1,2)$-gram TF-IDF features (\texttt{max\_features}\,=\,50{,}000) and
Logistic Regression with balanced class weights, the model attains
\textbf{0.431} accuracy and \textbf{0.424} macro-F1 on the validation set
(Table~\ref{tab:cls}). Although accuracy is close to the majority-class rate
($0.427$), the macro-F1 is far above the corresponding majority baseline
($\approx 0.20$), confirming that the model makes balanced predictions across
all three classes rather than collapsing to the majority. The confusion matrix
(Table~\ref{tab:cm}) shows the strongest confusion between Substitution and
Explanation, consistent with the fact that both often preserve sentence length.
These results establish a weak-but-nontrivial lower bound and motivate the
contextual models, which can capture semantic and structural cues unavailable
to bag-of-words features.

\begin{table}[h]
\centering
\begin{tabular}{lrrr}
\toprule
 & Precision & Recall & F1 \\
\midrule
Substitution   & 0.471 & 0.481 & 0.476 \\
Explanation    & 0.427 & 0.351 & 0.385 \\
Generalization & 0.376 & 0.449 & 0.410 \\
\midrule
Macro avg      & 0.425 & 0.427 & 0.424 \\
Accuracy       & \multicolumn{3}{c}{0.431 (n = 1{,}431)} \\
\bottomrule
\end{tabular}
\caption{TF-IDF + Logistic Regression, per-class results on the validation set.}
\label{tab:cls}
\end{table}

\begin{table}[h]
\centering
\begin{tabular}{lrrr}
\toprule
True $\backslash$ Pred & Sub & Exp & Gen \\
\midrule
Substitution   & 294 & 150 & 167 \\
Explanation    & 203 & 163 &  98 \\
Generalization & 127 &  69 & 160 \\
\bottomrule
\end{tabular}
\caption{Confusion matrix (rows = true, columns = predicted).}
\label{tab:cm}
\end{table}

%======================================================================
\section{Baselines}
%======================================================================

We frame our baselines as a comparison ladder anchored on \emph{warning
preservation} --- the fraction of source warning phrases that survive
simplification. This metric targets safety-critical information loss that
n-gram metrics such as ROUGE ignore.

\begin{table}[h]
\centering
\begin{tabular}{lc}
\toprule
Baseline & Warning preservation \\
\midrule
No simplification (ceiling)         & 1.000 \\
Rule-based CHV substitution         & 0.994 \\
Direct LLM prompting (FLAN-T5-Base) & 0.950 \\
Operation-aware system (ours)       & TBD \\
\bottomrule
\end{tabular}
\caption{Warning preservation across baselines. Direct LLM prompting loses the
most safety-critical information.}
\label{tab:baselines}
\end{table}

\textbf{Baseline 1 (no simplification)} leaves the text unchanged and thus
trivially preserves all warnings (1.000); it is the safety ceiling against which
no system should be expected to improve. \textbf{Baseline 2 (rule-based CHV
substitution)} performs Consumer Health Vocabulary term replacement and sentence
splitting, preserving 99.4\% of warnings --- auditable and low-risk, but limited
to vocabulary-level edits. \textbf{Baseline 3 (direct LLM prompting)} is the
deliberately naive approach our system argues against: we prompt an
instruction-tuned open-source LLM (FLAN-T5-Base) to ``simplify this biomedical
sentence for a patient,'' with no operation guidance and no safety constraint.
We run it locally to keep the baseline fully reproducible (no paid API).

On a 30-sentence sample restricted to warning-bearing sentences, direct LLM
prompting preserved only \textbf{95.0\%} of source warnings --- measurably below
both the rule-based baseline (99.4\%) and the do-nothing ceiling (100\%). This
confirms our central premise: unguided LLM simplification sacrifices
safety-critical information that conservative methods retain, motivating the
operation-aware routing our full system introduces. The three baselines together
bracket the safety--fluency trade-off our system must navigate.

%======================================================================
\section*{Literature Review}
%======================================================================

\textbf{SALSA edit taxonomy}~\cite{salsa} defines the fine-grained edit
operations that underpin our three-way operation taxonomy, directly grounding
our Substitution/Explanation/Generalization labels. The \textbf{PLABA
dataset}~\cite{plaba} provides the sentence-aligned source--adaptation pairs we
pseudo-label. \textbf{BERT}~\cite{bert} and \textbf{BioBERT}~\cite{biobert}
motivate the contextual and biomedical-domain classifiers in our model
progression. The \textbf{TSAR shared task}~\cite{tsar} documents the field's
shift toward LLM-based simplification and the accompanying safety gaps that
motivate our diagnostic, safety-first framing.

% VERIFY: the two entries below need confirmed bibliographic details before
% submission. Fill in authors/venue/year, then uncomment the \bibitem lines.
% - EhiMeNLP TSAR winner (best LLM ensemble)
% - Health QA SPQA (style perturbations reduce correctness/completeness)

\begin{thebibliography}{9}

\bibitem{salsa}
Heineman, D., Dou, Y., Maddela, M., \& Xu, W. (2023).
\emph{Dancing Between Success and Failure: Edit-level Simplification Evaluation
using SALSA.} EMNLP 2023. arXiv:2305.14458.

\bibitem{plaba}
Attal, K., Ondov, B., \& Demner-Fushman, D. (2023).
\emph{A dataset for plain language adaptation of biomedical abstracts.}
Scientific Data, 10(1), 8.

\bibitem{bert}
Devlin, J., Chang, M.-W., Lee, K., \& Toutanova, K. (2019).
\emph{BERT: Pre-training of Deep Bidirectional Transformers for Language
Understanding.} NAACL-HLT.

\bibitem{biobert}
Lee, J., Yoon, W., Kim, S., Kim, D., Kim, S., So, C.~H., \& Kang, J. (2020).
\emph{BioBERT: a pre-trained biomedical language representation model for
biomedical text mining.} Bioinformatics, 36(4), 1234--1240.

\bibitem{tsar}
Saggion, H., \v{S}tajner, S., Ferr\'es, D., et al. (2022).
\emph{Findings of the TSAR-2022 Shared Task on Multilingual Lexical
Simplification.} Proceedings of the TSAR Workshop, EMNLP 2022.
% VERIFY exact author list / page numbers.

% \bibitem{ehimenlp}  % VERIFY — TSAR winning LLM ensemble system
% \bibitem{spqa}      % VERIFY — Health QA SPQA style-perturbation paper
%========================= INTRODUCTION =========================
\section{Introduction}

Biomedical abstracts are written for clinicians, not patients. They are dense
with medical jargon, long and syntactically complex sentences, numerical risk
data, and safety warnings that a lay reader cannot follow. This creates a
health literacy gap: the people who most need to understand a treatment, a
risk, or a warning are often unable to read the text that describes it. Low
health literacy is linked to worse health outcomes, so making biomedical text
readable is a practical clinical concern, not only a linguistic one.

Most automatic simplification systems treat every kind of complexity the same
way. A sentence is fed into a model and a simpler sentence is returned, whether
the difficulty came from a single jargon term, a tangled clause structure, or a
critical dosage. Directly prompting a language model improves surface
readability, but it routinely drops medical warnings, weakens risk language, and
removes clinically important numbers. In other words, uniform simplification
buys readability at the cost of safety.

We take a diagnostic approach: classify first, then simplify. Before any text is
rewritten, we identify what type of complexity a span carries, and only then
apply the operation suited to it. A jargon term can be replaced from a
controlled vocabulary with no model involved. An unfamiliar concept can be
explained. An overly specific statement can be generalized. Routing each span to
the right operation lets us simplify aggressively where it is safe and stay
conservative where information must be preserved.

Our central research question is: \emph{does predicting the required
simplification operation at the span level reduce safety-critical information
loss, compared to rule-based, trained model, and direct LLM simplification
baselines?} We measure this with custom preservation metrics for entities,
warning phrases, and numerical expressions, alongside standard readability and
similarity measures.

%============================= DATA =============================
\section{Data}

We use PLABA, the Plain Language Adaptation of Biomedical Abstracts dataset
\cite{attal2023plaba}, released by Attal et al. and used as the official
benchmark at the TREC 2023 and 2024 PLABA tracks. PLABA contains 750 PubMed
abstracts retrieved to answer 75 consumer health questions, each manually
rewritten into plain language by at least one biomedical expert, for 921 gold
adaptations in total. The data is sentence-aligned, so every source sentence
maps to its simplified counterpart, and up to four reference simplifications are
available per sentence for SARI evaluation.

We work with a train, validation, and test split of 635, 138, and 148 sentence
pairs. An exploratory analysis of the 635 training sentences shows how tangled
the complexity is: 87.9\% contain medical jargon, 82.7\% are syntactically
complex beyond vocabulary-level fixes, 60.9\% carry numerical risk data such as
dosages or odds ratios, and 33.1\% contain safety-critical warning language.Biomedical named entity recognition flags diseases and chemicals requiring careful terminology handling in 87.9\% of sentences, often co-occurring with UMLS jargon signals.
Critically, 43.6\% of examples carry exactly three of these complexity types and
60.8\% carry three or more at once. A system that applies a single operation to
every span will therefore mishandle most inputs in some way, which is the core
motivation for our operation-aware design.

PLABA does not ship with the span-level operation labels our classifier needs,
so we derive them through pseudo-labeling from the aligned source and reference
sentences. We scope the labels to three operations, Substitution, Explanation,
and Generalization, which are the most frequent and most clinically distinct.
The Consumer Health Vocabulary and UMLS are used as external knowledge bases in
the pipeline rather than as training data.

%========================= RELATED WORK =========================
\section{Related Work}

\paragraph{Datasets.}
PLABA \cite{attal2023plaba} is the dataset our project is built on. It pairs 750
biomedical abstracts with expert plain-language rewrites at the sentence level
and annotates term-level edit types, which makes it well suited to studying
which operation a given span requires. Its sentence alignment is what lets us
build source-to-reference pairs and pseudo-label operations rather than treating
simplification as an opaque sequence-to-sequence task.

\paragraph{Prior PLABA systems.}
Knappich et al. \cite{knappich2023boschai} describe the BoschAI system that
ranked first in the PLABA 2023 shared task. Building on Llama 2, they observe
that source and target share most tokens, which gives a weak training signal and
produces models that edit conservatively. Their fix is to weight the loss toward
edited tokens using edit distance and edit operations. Their central finding,
that fine-tuned models tend to copy the input and edit little, directly
motivates our diagnostic framing: if a model barely edits, it also barely
distinguishes a jargon substitution from a critical warning that must be kept.
Li et al. \cite{li2024beemanc} present the BeeManc system for the same track,
investigating large language models with controllable attributes for biomedical
simplification, and provide a strong LLM-based point of comparison for our
operation-guided generation.

\paragraph{Evaluation metrics.}
SARI \cite{xu2016sari} is the primary simplification metric and our main quality
measure. Unlike BLEU or ROUGE, it compares the system output against both the
source and multiple references, rewarding words that are correctly added,
kept, and deleted, which is why PLABA provides several references per sentence.
BERTScore \cite{zhang2020bertscore} complements SARI by measuring semantic
similarity through contextual embeddings rather than surface overlap; we compute
it with a biomedical model so that domain terms are matched fairly. Neither
metric captures whether a warning or a dosage survived simplification, which is
the gap our custom preservation metrics are designed to fill.

\paragraph{Health literacy and readability.}
Berkman et al. \cite{berkman2011health} show through a systematic review that low
health literacy is associated with worse health outcomes, which grounds the
practical need for this work. Readability itself is commonly measured with the
Flesch-Kincaid Grade Level \cite{kincaid1975fkgl}, which we report so that
readability gains can be weighed against the safety costs they may carry.

%========================= REFERENCES =========================
\begin{thebibliography}{9}

\bibitem{attal2023plaba}
Attal, K., Ondov, B., \& Demner-Fushman, D. (2023).
A dataset for plain language adaptation of biomedical abstracts.
\emph{Scientific Data}, 10(1), 8.

\bibitem{knappich2023boschai}
Knappich, V., Razniewski, S., \& Friedrich, A. (2023).
BoschAI @ PLABA 2023: Leveraging edit operations in end-to-end neural sentence
simplification. \emph{arXiv preprint arXiv:2311.01907}.

\bibitem{li2024beemanc}
Li, Z., Belkadi, S., Micheletti, N., Han, L., Shardlow, M., \& Nenadic, G.
(2024). BeeManc at the PLABA track of TAC-2023: Investigating LLMs and
controllable attributes for improving biomedical text readability.
\emph{arXiv preprint}.

\bibitem{xu2016sari}
Xu, W., Napoles, C., Pavlick, E., Chen, Q., \& Callison-Burch, C. (2016).
Optimizing statistical machine translation for text simplification.
\emph{Transactions of the Association for Computational Linguistics}, 4,
401--415.

\bibitem{zhang2020bertscore}
Zhang, T., Kishore, V., Wu, F., Weinberger, K. Q., \& Artzi, Y. (2020).
BERTScore: Evaluating text generation with BERT.
\emph{International Conference on Learning Representations (ICLR)}.

\bibitem{berkman2011health}
Berkman, N. D., Sheridan, S. L., Donahue, K. E., Halpern, D. J., \& Crotty, K.
(2011). Low health literacy and health outcomes: An updated systematic review.
\emph{Annals of Internal Medicine}, 155(2), 97--107.

\bibitem{kincaid1975fkgl}
Kincaid, J. P., Fishburne, R. P., Rogers, R. L., \& Chissom, B. S. (1975).
Derivation of new readability formulas for Navy enlisted personnel.
\emph{Naval Technical Training Command Research Branch Report 8-75}.

\end{thebibliography}

\end{document}
